\documentclass[a4paper,12pt]{article} 

\usepackage[spanish]{babel}
\usepackage[utf8]{inputenc}
\usepackage[T1]{fontenc}

\usepackage{geometry}
\geometry{margin=2.5cm}

\usepackage{hyperref}
\usepackage{microtype} 

\usepackage{fancyhdr}

\pagestyle{fancy}
\fancyhf{}
\lhead{Class Notes}
\rhead{The Culture of Business}
\cfoot{\thepage}

\title{\vspace{-2cm}\textbf{Class Notes}\\[0.2cm]
\large The Culture of Business}
\author{Juan Ignacio Elosegui}
\date{Universidad Torcuato Di Tella \\ Marzo 2026}

\begin{document}
    \maketitle
    \section*{First Part}
        \subsection*{Class I - 03.03}
            \noindent
        \subsection*{Class II - 05.03}
            \noindent
        \subsection*{Class III - 10.03}
            \noindent
        \subsection*{Class IV - 12.03}
            \noindent
        \subsection*{Class V - 17.03}
            \noindent
        \subsection*{Class VI - 19.03}
            \noindent No economics if mathematics didn't exist. Economics wouldn't make sense. \\
            \\
            \noindent What does trade have to do with gravity? "to flow" and "to swell" also have to do with physics, some words used in businesses are also used in physics! Such as liquidity, flows, pipes, bubbles, and so on. There IS a reason for this. \\
            \\
            \noindent Mathematics is the language of physics, and physics have to do with economics as we can tell. There was an urge to create such a connection, it didn't happen naturally. Some concepts were borrowed from natural sciences. \\
            \\
            \noindent But why physics? Why did economics get such an unhealthy obsession with physics?Why not astronomy, or biology? Because physics was at the end of the 18th century the queen of sciences. It had plenty of very important discoveries. It was the closest we've ever felt to explain the theory of everything.

            Back in the day, Fisher used fluids as money. He gave flesh to his ideas of money with mere water and balances. Technically, he explained economics with physical phenomena with mechanisms.

            Bill Phillips used MONIAC: a hydraulic machine to represent stocks and flows of a model using water.

            Walras linked economics to Newton's laws: he believed economics was a mathematical science.

            We can also find a straight correlations between physics terms and economics terms: particle-individual, space-commodity, and so on. 

            This straight correlation is still among us today. \\
            \\
            \noindent Economics has gotten such authority because it could be explained with the queen of all sciences. Economics might make mistakes, but they can never be wrong; if you disagree with the theory of gravity you are a fool... but stock prices can be explained with this very same law, so it has the same authority and it is difficult to be challenged since you are not against an economist, but nature.
        \subsection*{Class VII}
            \noindent
        \subsection*{Class VIII}
            \noindent
\end{document}
